\fchapter{Ej 5 Pág 75}

\fsection{\boldit{\textcolor{waveBlue2}{Actividades}}}{ Encuentra los tres errores que aparecen en cada uno de los siguientes textos y redáctalos de manera correcta.}


\begin{solution}
\fsubsection{Texto 1}
Una nube pública es un conjunto de recursos proporcionados por una entidad externa, accesibles a cualquier usuario 
\textcolor{samuraiRed}{\sout{sin necesidad de utilizar internet}} \textcolor{springGreen}{con acceso a internet}. Ofrece escalabilidad bajo demanda, pago por uso y el proveedor se encarga del mantenimiento.

La nube privada, por otro lado, es exclusiva de una organización, lo que permite mayor control y \textcolor{samuraiRed}{\sout{ reduce los costes.}} \textcolor{springGreen}{aunque aumenta los costes a menor escala}.

La nube híbrida combina lo mejor de ambas, permitiendo flexibilidad y cumplimiento de regulaciones, mientras que el modelo multinube utiliza servicios  \textcolor{samuraiRed}{\sout{de un único proveedor}} \textcolor{springGreen}{de distintos proveedores} para maximizar beneficios y reducir dependencias.

	
\vspace{2em}
\hrulefill
\vspace{2em}

\fsubsection{Texto 2}
El modelo SaaS permite a los usuarios acceder a aplicaciones a través de internet sin necesidad de instalación local, con gestión \textcolor{samuraiRed}{\sout{centralizada por parte del cliente}} \textcolor{springGreen}{responsabilidad del proveedor} .

Se caracteriza por su escalabilidad y facturación por suscripción.

Existen soluciones SaaS tanto para empresas (B2B), como CRM y herramientas de marketing, como para consumidores (B2C), como servicios de streaming.

Las soluciones pueden ser verticales,  \textcolor{samuraiRed}{\sout{que cubren necesidades generales de negocios}}

\textcolor{springGreen}{las cuales ofrecen funcionalidades especializadas para sectores específicos}, u horizontales,\textcolor{samuraiRed}{\sout{las cuales ofrecen funcionalidades especializadas para sectores específicos}} 
\textcolor{springGreen}{que cubren necesidades generales de negocios}.


\end{solution}
